\section{Ordenamiento, Agrupación y Agregación}

\subsection{Ordenamiento}

Un \textit{Data Frame} puede estar desordenado tras una carga de datos o una limpieza masiva. \textit{Pandas} nos permite reorganizarlo tanto por sus etiquetas (\textit{índices}) como por sus valores reales.

\subsubsection{Ordenamiento por índice (\textit{sort\_index()})}

Si su \textit{Data Frame} viene desordenado, puede reorganizar por índice de forma ascendente con la función \texttt{.sort\_index()}.
\begin{pycode}{Ordenar por índice}
import pandas as pd

datos = {
    'Nombres': ['Juan', 'Pedro', 'Gabriela'],
    'Puestos': ['Administrativo', 'Secretario', 'Jefa'],
    'Salario': [200, 190, 250]
}

df = pd.DataFrame(datos, index=['z', 'v', 'm']) 
print(df) # Retorna:
#     Nombres         Puestos  Salario
# z      Juan  Administrativo      200
# v     Pedro      Secretario      190
# m  Gabriela            Jefa      250

df_ordenado = df.sort_index()
print(df_ordenado) # Retorna: 
#     Nombres         Puestos  Salario
# m  Gabriela            Jefa      250
# v     Pedro      Secretario      190
# z      Juan  Administrativo      200
\end{pycode}

Si desea ordenarlo de forma descendente utilice: 
\begin{verbatim}
    df.sort\_index(ascending=False)
\end{verbatim}


Para ordenar las \textit{columnas} de forma alfabética puede utilizar: \texttt{axis}, donde \texttt{axis=1} ordena las columnas.

\begin{pycode}{Ordenar columnas de forma descendente}
import pandas as pd

datos = {
    'Nombres': ['Juan', 'Pedro', 'Gabriela'],
    'Puestos': ['Administrativo', 'Secretario', 'Jefa'],
    'Salario': [200, 190, 250]
}

df = pd.DataFrame(datos, index=['z', 'v', 'm']) 
print(df) # Retorna:
#     Nombres         Puestos  Salario
# z      Juan  Administrativo      200
# v     Pedro      Secretario      190
# m  Gabriela            Jefa      250

df_columnas_ordenadas_descendentes = df.sort_index(axis=1, ascending=False)
print(df_columnas_ordenadas_descendentes) # Retorna: 
#    Salario         Puestos   Nombres
# z      200  Administrativo      Juan
# v      190      Secretario     Pedro
# m      250            Jefa  Gabriela
\end{pycode}

\begin{nota}
    Si no quiere crear una variable nueva (\texttt{df\_ordenado}), puede usar el parámetro \texttt{inplace=True}. Esto modificará el \textit{Data Frame} original directamente: \texttt{df.sort\_index(inplace=True)}.
\end{nota}


\subsubsection{Ordenamiento por valores (\textit{sort\_values()})}

A diferencia del ordenamiento por índice, \texttt{.sort\_values()} reorganiza la tabla basándose en el contenido de una o más columnas específicas. Al utilizar la función \texttt{.sort\_values} ordenamos los valores de la columna de forma ascendente.

Para ejecutarlo, es obligatorio indicar qué columna servirá como guía mediante el parámetro:
\begin{verbatim}
    df.sort_values(by='columna')
\end{verbatim}

Donde la columna va a ser la que va a estar organizada de forma ascendente.

\begin{pycode}{Ordenar valores de forma ascendente}
import pandas as pd

datos = {
    'Nombres': ['Juan', 'Pedro', 'Gabriela'],
    'Puestos': ['Administrativo', 'Secretario', 'Jefa'],
    'Salario': [200, 190, 250]
}

df = pd.DataFrame(datos, index=['z', 'v', 'm']) 
print(df) # Retorna:
#     Nombres         Puestos  Salario
# z      Juan  Administrativo      200
# v     Pedro      Secretario      190
# m  Gabriela            Jefa      250

df.sort_values(by='Salario', inplace=True)
print(df) # Return 
#     Nombres         Puestos  Salario
# v     Pedro      Secretario      190
# z      Juan  Administrativo      200
# m  Gabriela            Jefa      250
\end{pycode}

Al igual que en \texttt{sort\_index()} se puede ordenar de forma descendente con \texttt{ascending=False}
\begin{pycode}{Valores por orden descendente}
import pandas as pd

datos = {
    'Nombres': ['Juan', 'Pedro', 'Gabriela'],
    'Puestos': ['Administrativo', 'Secretario', 'Jefa'],
    'Salario': [200, 190, 250]
}

df = pd.DataFrame(datos, index=['z', 'v', 'm']) 
print(df) # Retorna:
#     Nombres         Puestos  Salario
# z      Juan  Administrativo      200
# v     Pedro      Secretario      190
# m  Gabriela            Jefa      250

df.sort_values(by='Puestos', ascending=False, inplace=True)
print(df) # Return 
#     Nombres         Puestos  Salario
# v     Pedro      Secretario      190
# m  Gabriela            Jefa      250
# z      Juan  Administrativo      200
\end{pycode}

\begin{nota}
    Recuerde que las letras tambien son valores, se ordenan alfabéticamente.
\end{nota}

También se puede ordenar por varias columnas. Primero se ordena por la primera columna, y en caso de que se repita se ordena por la segunda columna que se otorgo. Esto se veria así:
\begin{verbatim}
    df.sort_values(by=["Puestos", "Salario"])
\end{verbatim}
Donde se va a ordenar por \texttt{Puestos}, pero si dos puestos son iguales, se va a ordenar por \texttt{Salario}.

\begin{nota}
    Si quiere que una columna sea ascendente y la otra descendente, puede pasar una lista al parámetro \texttt{ascending}: \texttt{df.sort\_values(by=['Puestos', 'Salario'], ascending=[True, False])}
\end{nota}

\begin{nota}
    En el caso de tener valores iguales, se puede utilizar \texttt{kind='mergesort'} dentro de \texttt{sort\_values()}. Para aquellos valores repitidos mantengan el orden original que tenian en la lista vieja.
\end{nota}



\subsection{Agregación simple}

Las funciones de agregación son aquellas que permiten obtener indicadores. Estas son:
\begin{verbatim}
    .mean() 
    .sum()
    .count()
\end{verbatim}



\subsection{Agrupación (\textit{.groupby()})}

La función \texttt{.groupby} permite segmentar los datos para realizar cálculos por categorías.

Su proceso consiste en:
\begin{enumerate}
    \item \textit{Split} (Dividir): Separamos los datos en grupos basados en columnas elegidas con \texttt{.groupby()}.
    \item \textit{Apply} (Aplicar): Ejecutamos una operación sobre cada grupo de fomra independiente.
    \item \textit{Combine} (Combinar): Juntamos los resultados de esos calculos en una nueva estructura de datos.
\end{enumerate}

\begin{pycode}{Utilización de .groupby()}
import pandas as pd

datos = {
    'Programa': ['python', 'c', 'python', 'java', 'c', 'python'],
    'Salario': [100, 50, 120, 80, 60, 90]
}
df = pd.DataFrame(datos)

# Agrupamos por 'Programa' y calculamos la media del 'Salario'
resultado = df.groupby('Programa', as_index=False)['Salario'].mean()

print(resultado) # Retorna:
#   Programa     Salario
# 0        c   55.000000
# 1     java   80.000000
# 2   python  103.333333
\end{pycode}

\begin{nota}
    Luego de \texttt{.groupby()}, \textit{Pandas} convierte las columnas por las que agrupas en el Índice del nuevo DataFrame. Para que siga siendo una columna normal, es decir que aparezca la columna salario, utilizamos \texttt{as\_index=False}, que básicamente luego de agrupar los datos, ordena que las columnas se mantengan como las originales.
\end{nota}

También se puede agrupar creando una lista de columnas. Por ejemplo:
\begin{verbatim}
    df.groupby(['columna1', 'columna2'])['columna'].mean()
\end{verbatim}
Donde se crean subgrupos dentro de grupos. 
\begin{pycode}{Agrupación múltiple}
import pandas as pd

datos = {
    'Sede': ['Madrid', 'Madrid', 'Barcelona', 'Barcelona', 'Madrid'],
    'Lenguaje': ['Python', 'Java', 'Python', 'Python', 'Java'],
    'Salario': [100, 90, 110, 105, 95]
}
df = pd.DataFrame(datos)

# Agrupamos por Sede Y LENGUAJE
agrupado = df.groupby(['Sede', 'Lenguaje'])['Salario'].mean()

print(agrupado) # Retorna:
# Sede       Lenguaje
# Barcelona  Python      107.5
# Madrid     Java         92.5
#            Python      100.0

# Madrid es un grupo y dentro tiene el subgurpo de java y python.
\end{pycode}




\subsection{Agregaciones múltiples (\textit{.agg()})}

La función \texttt{.agg()}(aggregate), significa tomar más de un valor y resumirlo en uno solo. Esta combina las funciones básicas en una sola llamada.

\subsubsection{Agregación múltiple}

Para agregar más de una funcion múltiple podemos utilizar \texttt{.agg(funciones)}
\begin{pycode}{Agregación Múltiple .agg()}
import pandas as pd

df = pd.DataFrame({
    'Programa': ['Python', 'Python', 'Java', 'Java'],
    'Salario': [100, 150, 80, 90]
})

# Pedimos varias funciones a la vez pasando una LISTA []
resumen = df.groupby('Programa')['Salario'].agg(['min', 'max', 'sum', 'mean'])

print(resumen) # Retorna:
#           min  max  sum   mean
# Programa                      
# Java       80   90  170   85.0
# Python    100  150  250  125.0
\end{pycode}

\subsubsection{Agregación por diccionario}

Si hay un \textit{Data Frame} con múltiples columnas, podemos definir que calculo se quiere para cada una mediante el uso de un \textit{diccionario}.
\begin{pycode}{Agrupación por diccionario}
import pandas as pd

datos = {
    'Departamento': ['I+D', 'I+D', 'Ventas', 'Ventas'],
    'Salario': [100, 150, 80, 90],
    'Experiencia': [5, 8, 2, 3]
}
df = pd.DataFrame(datos)

# Para Salario queremos el promedio, para Experiencia el máximo
reglas = {
    'Salario': 'mean',
    'Experiencia': 'max'
}

resultado = df.groupby('Departamento').agg(reglas)
print(resultado) # Retorna:
#               Salario  Experiencia
# Departamento                      
# I+D             125.0            8
# Ventas           85.0            3    
\end{pycode}

\subsubsection{Agregación con nombre personalizado}

A las funciones se le puede colocar el nombre en la misma funcion \texttt{.agg()}, donde:
\begin{verbatim}
    .agg(nombre_nuevo = ('columna', 'función'))
\end{verbatim}

\begin{pycode}{Agregación con nombre personalizado}
import pandas as pd

datos = {
    'Departamento': ['I+D', 'I+D', 'Ventas', 'Ventas'],
    'Salario': [100, 150, 80, 90],
    'Experiencia': [5, 8, 2, 3]
}
df = pd.DataFrame(datos)

# Sintaxis: nombre_nuevo = ('columna_original', 'funcion')
resumen = df.groupby('Departamento').agg(
    salario_promedio = ('Salario', 'mean'),
    max_experiencia = ('Experiencia', 'max')
)

print(resumen) # Retorna:
#               salario_promedio  max_experiencia
# Departamento                                   
# I+D                      125.0                8
# Ventas                    85.0                3 
\end{pycode}

\begin{nota}
    Otras funciones que representan operaciones matemáticas son: \texttt{'sum', 'mean', 'std', 'count', 'min', 'max'}. Tambien pueden haber funciones de librerias de \textit{NumPy} o funciones personalizadas.
\end{nota}





\subsection{Tansformar datos (\textit{.transform()})}

La función de \texttt{.transform()} es utilizada para realizar un cálculo por grupo, pero manteniendo el mismo número de filas que el \textit{Data Frame} inicial.

Si usted quisiera comparar el sueldo de cada empleado contra el sueldo de su departamento, con \texttt{transform()} puede calcular el promnedio de cada grupo y este lo ``reparte'' en todas las filas pertenecientes a ese grupo.

\begin{pycode}{.transform()}
import pandas as pd

df = pd.DataFrame({
    'Nombre': ['Juan', 'Ana', 'Pedro', 'Laura'],
    'Lenguaje': ['Python', 'Python', 'Java', 'Java'],
    'Salario': [100, 150, 80, 90]
})

# Usamos transform para que el promedio se repita en cada fila
df['Promedio_Lenguaje'] = df.groupby('Lenguaje')['Salario'].transform('mean')

print(df) # Retorna:
#   Nombre Lenguaje  Salario  Promedio_Lenguaje
# 0   Juan   Python      100              125.0
# 1    Ana   Python      150              125.0
# 2  Pedro     Java       80               85.0
# 3  Laura     Java       90               85.0    
\end{pycode}


